

\subsection{Further related work}
Our work builds on prior research in multi-winner elections, metric distortion, and (fair) clustering. As these areas have been extensively studied, we only discuss works that are most related to our 
methodological approach.


Besides extending the framework of metric distortion and relating the problem with clustering~\cite{burkhardt2024low, pulyassary2025constant, caragiannis2022metric}, 
other lines of works include the study of the \emph{(approximately) stable committee selection} problem~\cite{fain2018fair,jiang2020approximately}, and selecting a committee of minimum size that is a \emph{Condorcet winner}~\cite{charikar2025six}.
Concurrently, there is a work which studies a \emph{peer selection} problem, where the voters and candidates coincide~\cite{cembrano2025metric}. We note from the technical perspective, they focus on the line metric without distance queries, and their techniques may not directly extend to our setting.


Clustering is a fundamental unsupervised machine learning task that has been extensively studied~\cite{jain1999data}. In recent years, awareness about automated decision-making propagating biases has led to an increase in attention towards algorithmic fairness principles. Consequently, several classic unsupervised learning problems, including clustering, have been reintroduced with fairness constraints~\cite{thejaswi2021diversity, thejaswi2022clustering,thejaswi2024diversity,gadekar2025capacitated,matakos2024fair,zhang2024parameterized,abbasi2023parameterized,gadekar2025fair,chierichetti2017fair}. Among these formulations, the one most relevant to our work is the fair $k$-center problem introduced by \citet{kleindessner2019fair}, which seeks to choose a specified number of representatives from each demographic group while minimizing the egalitarian (min-max) objective, \ie, minimizing the maximum distance from any client to its nearest representative. \citet{jones2020fair} presented a $3$-approximation algorithm using a matching framework. This framework has been extended to solve several fair clustering variants~\cite{gadekar2025fair,chen2024approximation}.

The remainder of the paper is organized as follows. 
Section~\ref{sec:problem-formulation} introduces the necessary terminology and problem definitions. 
Section~\ref{sec:conventional-algorithms} provides an overview of the algorithmic techniques that form the basis for our algorithmic methods and proposes a $3$-distortion algorithm with $\bigO(k^2)$ queries.
Section~\ref{sec:analysis-ordinal-fair-k-center-with-limited-distance-queries} presents our main algorithmic contribution, a $5$-distortion algorithm with $\bigO(k\log^2 k)$ queries.


 






